\documentclass[a4paper,french,8pt]{article}
\usepackage{babel,amssymb,amsmath}
\usepackage[utf8]{inputenc}
\usepackage{array,colortbl}
\usepackage[T1]{fontenc}
\usepackage{multicol}
%\usepackage{multirow}
%\usepackage[8pt]{pstricks}
\usepackage{textcomp}
\usepackage{fancyhdr}
\usepackage{geometry}
\usepackage{eurosym}
% \usepackage{titling}
%\usepackage{makeidx}
% \usepackage[maxfloats=50]{morefloats}
% \usepackage{hyperref}
\usepackage{ulem}
\usepackage{epigraph}
\geometry{rmargin=2.5cm, lmargin=2.5cm, hmargin=2cm, bmargin=2cm}
\pagestyle{plain}
\usepackage[pdftex]{thumbpdf}
%\definecolor{vert}{rgb}{0.25,0.65,0.04}
\usepackage[pdftex,
    bookmarks         = true,
    bookmarksnumbered = true,
    pdfpagemode       = None,
    pdfstartview      = FitH,
    pdfpagelayout     = SinglePage,
    colorlinks        = true,
    linkcolor	      = black,
    pdfborder         = {0 0 0}
    ]{hyperref}
\usepackage[pdftex]{graphicx}


\newcommand{\HRule}{\rule{\linewidth}{0.5mm}}
% \hypersetup{
% 	pdfauthor = {Philippe \textsc{GAULTIER}},
% 	pdftitle = {Rapport de stage du 17/06/2013 au 23/08/2013},
% 	pdfsubject = {},
% 	pdfkeywords = {},
% 	pdfcreator = {},
% }
% \newcommand{\filename}[1]{\textbf{\#1}}
% \newcommand{\fonction}[1]{\textbf{\itshape{\#1}}} 
% \newcommand{\tabl}[1]{\textbf{\itshape{\#1}}[ ]} 
\begin{document}
\begin{titlepage}
\begin{center}

%Logo
\includegraphics[width=0.4\textwidth]{./logo_ENSIIE.png}~\\[1cm]
\textsc{\huge Ensiie Strasbourg}\\[1.5cm]

%Title
\HRule \\[0.4cm]
{
	\huge \bfseries Rapport de stage
\\[0.4cm] }
\HRule \\[1.5cm]

% Author and supervisor
\begin{minipage}{0.4\textwidth}
\begin{flushleft} \huge
\emph{Auteur:}\\
Philippe \textsc{Gaultier},\\[0.5cm]
\Large élève ingénieur à l'Ensiie Strasbourg
\end{flushleft}
\end{minipage}
\begin{minipage}{0.4\textwidth}
\begin{flushright} \huge
\emph{Maître de stage:} \\
André \textsc{Schaaf},\\[0.5cm]
\Large enseignant-chercheur à l'Université de Strasbourg
\end{flushright}
\end{minipage}

\vfill

% Bottom of the page
{\large Strasbourg, le \today}



%\maketitle
% \theauthor
% \thetitle
% \thedate
% \makeindex
\end{center}
\end{titlepage}

\newpage
{
  \centering
  {
    \vspace{3cm}
    \epigraph{Always code as if the guy who ends up maintaining your code will be a violent psychopath who knows where you live}{Martin Golding}
    \vspace{3cm}
    \epigraph{The road is long and in the end the journey is the destination}{Unknown}
  }
}
\newpage
\textit{\normalsize Dans toute la suite du rapport, l'<<Observatoire>> désigne l'<<Observatoire astronomique de Strasbourg>>
et l'<<Unistra>> ou l'<<UDS>> désignent l'<<Université de Strasbourg>>.}
\setlength{\columnseprule}{0.5pt}
\tableofcontents

\newpage

\section{Introduction}

	L'Observatoire est un établissement de recherche et d'enseignement centré sur l'astronomie, mais c'est aussi un 
	centre de données astronomiques et un centre d'observation réputés mondialement. Il représente la continuité entre
	l'ancien et le nouveau car il dispose d'un riche patrimoine mais est aussi à la pointe de la recherche.  \\
	De plus l'Observatoire est le parfait exemple de l'informatique au service d'autres spécialités scientifiques,
	à la fois dans le domaine de l'expertise, mais aussi sous un aspect éducatif.\\
	Pour toutes ces raisons, j'ai choisi d'effectuer mon stage de deuxième année au sein de l'Observatoire,
	au contact de technologies émergentes à savoir l'Oculus Rift et le rendu graphique 3D moderne.

\section{Présentation de l'observatoire}

	\subsection{Histoire}
		L'Observatoire a été fondé en 1881 sur l'initiative de l'empereur Guillaume II, l'Alsace étant allemande
		à cette époque.\\
		Il est constitué de trois bâtiments : une Grande Coupole, un bâtiment des salles méridiennes avec deux coupoles,
		et un bâtiment à usage de bureau et de résidence.\\
		La Grande Coupole en fer, de 9,2 mètres de diamètre et pesant 34 tonnes2, contient le Grand Réfracteur,
		une lunette de 48,7 cm d'ouverture et 7 m de focale, construite en 1877, la plus grande d'Europe
		au moment de son installation et aujourd'hui (2008) la troisième de France en taille.\\
		Il dispose également d’un riche patrimoine d’instruments et d’ouvrage anciens.
		
	\subsection{Centre de données astronomiques de Strasbourg (CDS)}
	
		Le CDS est à la fois une équipe de recherche et un Service d’Observation.
		Les services de bases de données (SIMBAD, Vizier) et de visualisation (ALADIN) développés par le CDS
		sont utilisés par l’ensemble de la communauté astronomique mondiale. \\
		Celui-ci est l'un des acteurs majeurs du développement de l'Observatoire Virtuel International en astronomie.
		Fin 2008, le CDS a été labellisé TGIR (Très Grande Infrastructure de Recherche) par le Ministère de l'Enseignement Supérieur et de la Recherche, 
		reconfirmé comme Infrastructure de Recherche en 2012, ce qui le range au même niveau que des infrastructures internationales
		comme l’European Southern Observatory ou RENATER à l’échelon national.

	\subsection{Equipe de recherche Galaxies}
	
		L’équipe <<Galaxies>> étudie la formation et l’évolution des galaxies et de notre Galaxie 
		au travers de leurs populations stellaires et de la dynamique des étoiles et de la matière noire. \\
		Elle est impliquée dans la préparation de la mission satellitaire astrométrique Gaia de l’Agence Spatiale Européenne 
		dont le lancement est prévu en 2012 et dans le grand relevé cinématique RAVE. 
	
	\subsection{Equipe de recherche Hautes Énergies} 
	
		L’équipe <<Hautes Énergies>> s’intéresse aux sources galactiques et extragalactiques émettrices en rayons X,
		objets compacts (étoiles à neutron, naines blanches, etc.) et noyaux actifs de galaxies.\\
		Elle est impliquée dans le SSC-XMM, un consortium international de laboratoires sélectionné par l’ESA
		et labellisé par l’INSU comme Service d’Observation, qui est en charge de fournir des catalogues complets
		de sources X observées par le satellite XMM-Newton à la communauté internationale. 

\section{Mon stage}

	\subsection{Objectif}	
	
		L’objectif de ce stage est centré autour de l'Oculus Rift et est double:
		\begin{itemize}
		 \item Intégration de l'Oculus Rift à une simulation 3D du système solaire existante
		 \item Développement d'un programme de visualisation 3D  d'étoiles avec intégration de l'Oculus Rift
		\end{itemize}

	
	\subsection{L'Oculus Rift}
		
		L'Oculus Rift est un masque de réalité virtuelle, développé par \emph{Oculus VR}, une entreprise 
		basée en Californie et rachetée par \emph{Facebook} en mars 2014 pour  2 milliards \$. \\
		L'Oculus Rift a été initialement financé via une plateforme de financement collaboratif, \emph{Kickstarter},
		et a levé 91 millions \$ à cette occasion. \\
		Il permet une immersion  réaliste dans une scène en trois dimensions, en donnant l'impression d'y être physiquement
		présent, et crée ainsi une nouvelle expérience
		utilisateur. \\
		De plus, son prix est relativement peu élevé (350 \$, environ 300 \euro), ce qui le rend accessible au grand public. \\
		Pour toutes ces raisons, l'Oculus Rift est adapté à un usage éducatif et professionel, dans des domaines aussi variés
		que la simulation scientifique, le divertissement, l'éducation, \ldots.
		
		La version grand public est prévue pour fin 2014 ou début 2015. J'ai pour ma part travaillé avec la première
		version du masque, le \emph{DK1}, tandis que la deuxième version, le \emph{DK2} a été distribuée à partir d'août 2014.
		

	\subsection{Les nouvelles fonctionnalités}
	
		Les nouvelles fonctionnalités à implémenter sont diverses:
		
		\begin{itemize}
			\item Mise en place de groupes d’habilitation pour restreindre les accès et les droits sur les données (2 types d'accès: visualisation seule ou accès total)
			\item Améliorer le confort utilisateur en proposant des critères de recherche sur les données
			\item Implémenter l’activation/désactivation d’un compte (qui se faisait manuellement auparavant via une sauvegarde dans un fichier Excel externe à l'application)
			\item Empêcher que certaines données (notamment les cumuls) soient modifiées
		\end{itemize}
		
	\subsection{Contraintes}
		
		L'application web doit respecter les normes de programmation et de qualité du CM-CIC. Plus précisément, certains mot-clés et certaines fonctionnalités des langages sont interdits d'utilisation,
		notamment:
		
		\begin{itemize}
			\item L'usage des pointeurs en C\# et en COBOL
			\item L'utilisation de Javascript dans l'interface web
			\item Le stockage de variables générales "en dur" dans le code (comme les tables de correspondances code société-libellé)
			\item L'utilisation certains mot-clés comme "STOP RUN", "EXIT PROGRAM", "EJECT", etc. en COBOL
		\end{itemize}
		
		De plus, le modèle d'architecture d'un projet est imposé (découpe en programmes métiers/logiques/physiques).
	
	\subsection{La norme ISO9001}
	
		La norme ISO9001 est un standard international  de qualité dont les exigences sont publiées par l’Organisation Internationale de Normalisation.
		Elle est validée par des sociétés d’audit et de certification externes. 
		A EI, tous les produits sont certifiés ISO9001. Les équipes subissent des audits internes et externes réguliers pour garantir cette norme. 
		Les contraintes qui en découlent sont une production de documentation accrue et un reporting permanent dans l’avancement du projet et la résolution des incidents.
		Mon projet est ainsi conforme à la norme ISO9001.

	\subsection{Outils utilisés}
	
		\subsubsection{Programmes utilisés}
		
			\begin{itemize}
				\item Microsoft Visual Source Safe 2005(gestionnaire de versions)
				\item Microsoft Visual Studio 2010 (développement C\#, Xaml et Xml)
				\item Microsoft Visio 2010 (schémas et analyse)
				\item IBM Rational Developer for System z (développement COBOL)
			\end{itemize}
		
		\subsubsection{Langages utilisés}
			
			\begin{itemize}
				\item C\# 4.0 (développement web)
				\item Xaml, HTML (Interface Homme-Machine)
				\item Xml (mise en forme des données)
				\item COBOL version 3 release 4 (traitements métiers)
				\item SQL DB2 (base de données)
			\end{itemize}
			
			Le Xaml m’a permis de mettre en forme les vues de l’application web : agencer les éléments entre eux, définir leurs propriétés,
			et même mettre en place un rafraichissement partiel de la page (Ajax) sans écrire de JavaScript. C’est la partie Vue.\\
			J’ai utilisé le Xml pour décrire les données qui alimentaient ces vues : c’est (partiellement) la partie Modèle. L’autre partie est la base de données DB2 (back-end).\\
			La partie C\# quant à elle a servi à faire le routage entre les différentes pages ainsi qu’à introduire un minimum de logique en faisant des vérifications sur les données entrées par l'utilisateur 
			et en faisant les appels aux programmes COBOL. C’est la partie Contrôleur.\\ 
						
			Ce sont les programmes COBOL qui implémentent la véritable logique de l’application avec les règles métier.
			C’est aussi par l’intermédiaire du COBOL que se font les appels SQL à la base DB2.
			L’interface entre la partie C\# et la partie COBOL se fait au moyen du format d’échange propre à EI : le LTNV. 
			Ce dernier est ensuite transformé par l’intermédiaire du C\# en Xml nécessaire à l’alimentation des vues de l’application web.\\ 
			
			Les données suivent ainsi le cheminement suivant : \\

			 Données SQL DB2 $\to$ Variables COBOL $\to$ LTNV $\to$ Variable C\# $\to$ Xml $\to$ Xaml $\to$ Utilisateur 


	
		\subsubsection{Frameworks utilisés}
		
			\begin{itemize}
				\item .Net 4.0 (framework microsoft pour le développement web)
				\item Devbooster (framework interne pour le développement web)
			\end{itemize}
			
			Toute la partie front-end, agencée selon le modèle Modèle Vue Controlleur, a été développée à la fois au sein du Framework  web propre à EI (Devbooster) et au sein du Framework .Net.
		
		
	\subsection{Déroulement du stage}
	
		La première semaine (17/06 - 21/06) a été dédiée à l’analyse du programme existant, à l’architecture de l’application web et au dialogue avec les utilisateurs,
		qui s’est poursuivi tout au long du stage. Elle a débuté par une rencontre avec les utilisateurs (17/06), où ont été présentés les besoins, les qualités et défauts de l’existant et les améliorations possibles. 
		Cette semaine s’est ensuite poursuivie avec le développement de la maquette de l'interface web. \\ 
		
		La deuxième semaine (24/06 – 28/06) a été centrée sur l’interface (IHM) et l’ergonomie, avec la poursuite de la réalisation de la maquette de l’application web.
		Elle a commencée avec un point sur l’interface homme-machine (24/06), l’ergonomie et les normes d’homogénéisation du groupe en termes d’interface.\\ 
		
		La troisième semaine (01/07 – 05/07) a prolongé la deuxième avec la transformation de la maquette en véritable application web, fonctionnant avec des données factices, 
		même si les traitements métiers n'étaient pas implémentés.
		Elle a été notamment l’occasion de faire un point sur l’avancement de l’interface web (03/07).\\ 

		La quatrième semaine (08/07 – 12/07) a marqué un cap puisque je me suis focalisé sur la partie métier de l’application (gérée en COBOL) avec l’apprentissage de ce langage. 
		En parallèle ont eu lieu deux réunions : une réunion d’équipe pour présenter et récapituler le travail de chacun (09/07) et une réunion avec les utilisateurs pour présenter la nouvelle interface (10/07). 
		Cette dernière réunion a été l’occasion de soulever de nombreuses modifications et fonctionnalités à apporter, tant en terme d’ergonomie que de métier.\\

		La cinquième semaine (15/07 – 19/07) a été consacrée à l’architecture et à l’analyse de la partie COBOL. Elle s’est achevée par une réunion d’équipe (19/07). \\ 
		
		La sixième semaine (22/07 - 26/07) a poursuivi l’architecture et à l’analyse de la partie COBOL, tout en peaufinant l'interface web en échangeant avec les utilisateurs sur de nombreux points.
		Un point technique COBOL a été fait le 22/07.\\ 
		
		la septième semaine (29/07 - 02/08) a consisté principalement en l'amélioration de l'ergonomie de l'interface web et à la préparation des appels depuis cette dernière à la partie COBOL. De plus,
		elle a été consacrée à la rédaction du rapport de stage.
		\\ 
		
		La huitième semaine (05/08 - 09/08), ainsi que la neuvième (12/08 - 16/08) et la dixième (19/08 - 23/08), ont été consacrées au développements COBOL 
		et à la mise en place des appels host depuis le web.\\ 
	
	\subsection{Bonnes pratiques de développement}
			Pour tout le développement de l'application web, j'ai suivi les conventions de nommage préconisées par Microsoft dans le document
			"C\# Coding Conventions (C\# Programming Guide)" (\url{http://msdn.microsoft.com/en-us/library/vstudio/ff926074.aspx}), c'est-à-dire:
			\begin{description} 
				\item[Conventions de nommage] \hfill \\
					Toutes les variables, noms de méthodes, noms de classes, etc. sont écrites avec la casse $Pascal$ (exemple: $MaMethode()$ ), sauf les membres privés et les paramètres de méthodes qui suivent
					la casse $Camelcase$.
					L'indentation, le style de code et la rédaction des commentaires sont aussi régis par des règles précises.
				\item[Refactoring] \hfill \\
					Vers la fin de mon projet, après une pause dans le développement de l'interface web pour me consacrer au développement COBOL, il m'est apparu que le code de l'interface web était moyennement clair pour moi, 
					et peu clair pour les autres, surtout à cause de noms de variables peu explicites et des corps de méthodes trop longues.
					J'ai donc pu faire la connaissance des outils de $refactoring$ de Visual Studio, en particulier $Refactor \to Rename$ pour renommer intelligiblement les variables
					et $Refactor \to Extract Method$ pour découper les méthodes trop longues en plusieurs méthodes courtes.
				\item[Test unitaires]  \hfill \\
					Je me suis heurté à un cas où la vérification des données saisies par l'utilisateur était faite incorrectement par le framework. J'ai donc développé mes propres méthodes pour valider les données saisies,
					notamment pour vérifier qu'un coefficient était du bon type et dans le bon intervalle. Me heurtant à certains bugs, j'ai écrit des test unitaires avec l'outil intégré de Visual Studio 
					pour valider ces méthodes. Ainsi, le risque de données incorrectes validées est diminué, et il est possible de faire des tests de non-régression si changements il devait y avoir.
					J'ai pu également ainsi automatiser ces tests au lieu de les faire manuellement au travers de l'interface, ce qui m'a permis de tester plus de cas et d'isoler ces tests de l'interface.
			\end{description}
	
	\subsection{Architecture}
	
		\subsubsection{Interface web}
		
			\begin{figure}[h!]
			\caption{Fichiers de l'interface web}
			\centering
%				\includegraphics[width=6cm]{fichiers.png} 
			\end{figure} 
			
			%%%%%%%%%%%%%%%%%%%%%%%%%%%%%%%%%%%%%%%%%
				\paragraph{Fichiers C\#}
				
					\begin{figure}[h!]
					  \caption{Fichiers C\# de l'interface web}
					  \centering
%						  \includegraphics[width=6cm]{fichiers_c.png} 
					\end{figure} 
					
					\begin{itemize}
						\item App\_Code/BusinessLogic/Compte.cs : Fichier décrivant la classe Compte
						\item App\_Code/BusinessLogic/Liaison.cs : Fichier décrivant la classe Liaison
						\item App\_Code/Commands/\_EntryPoint.cs : Procédures de vérification des ressources et des droits d’accès à l’application
						\item App\_Code/Commands/Menu.cs	: Code C\# associé à la page décrivant les entrées dans le menu latéral
						\item App\_Code/Commands/\* : Fichiers C\# associées aux feuilles Xaml, contenant le code de navigation et les appels aux programmes métier
					\end{itemize}

			%%%%%%%%%%%%%%%%%%%%%%%%%%%%%%%%%%%%%%%%%
				\paragraph{Fichiers Xaml}
				      \begin{figure}[h!]
					\caption{Fichiers Xaml de l'interface web}
					\centering
%						\includegraphics[width=6cm]{fichiers_xaml.png} 
					\end{figure} 
					
					\begin{itemize}
						\item ListeLiaison.xaml : Ecran de liste des liaisons louées
						\item DetailLiaison.xaml : Ecran de détail d’une liaison louée
						\item FormulaireLiaison.xaml : Ecran de modification d'une liaison
						%\item SupprimerLiaison.xaml : Ecran de confirmation de suppression d’une liaison
						\item ListeComptes.xaml	: Ecran de liste des comptes
						\item FormulaireCompte.xaml	: Ecran de création et de modification de compte
						%\item SupprimerCompte.xaml : Ecran de confirmation de suppression d’un compte
						\item Facture.xaml : Ecran récapitulatif de la facture et de confirmation de création
						\item Menu.xaml	: Le menu latéral avec les liens vers ListeLiaison.xaml, ListeComptes.xaml et Facture.xaml
					\end{itemize}

			%%%%%%%%%%%%%%%%%%%%%%%%%%%%%%%%%%%%%%%%%
				\paragraph{Fichiers Xml}
				      \begin{figure}[h!]
					\caption{Fichiers Xml de l'interface web}
					\centering
%						\includegraphics[width=6cm]{fichiers_xml.png} 
					\end{figure} 
					Aux pages Xaml et C\# sont associés les fichiers Xml qui décrivent les données: \\
					App\_GlobalResources/Model/* : Pages Xml associées aux pages Xaml et C\#, contenant les données.

			%%%%%%%%%%%%%%%%%%%%%%%%%%%%%%%%%%%%%%%%%				
				\paragraph{Fichiers de configuration}
% 				    \begin{figure}[h!]
% 					\caption{Fichiers de configuration de l'interface web}
% 					\centering
% 						\includegraphics[width=12cm]{fichiers_cfg_projet.png} 
% 					\end{figure} 
% 					
					Les fichiers suivants sont utilisés pour paramétrer l’application Web :
					\begin{itemize}
						\item /devbappli.config : Fichier contenant diverses options dont le flag "mode design" (true pour mettre en mode maquette)
						\item /web.config : Fichier contenant diverses options et le flag de compilation en mode debug
						\item /zz/devbtables.config	: Permet de stocker des tables de correspondances (par exemple code société-libellé)
					\end{itemize}
		
			%%%%%%%%%%%%%%%%%%%%%%%%%%%%%%%%%%%%%%%%%
			\begin{figure}[h!]
			\caption{Cinématique de l'interface web}
			\centering
%				\includegraphics[width=12cm]{cinematique_facll.png} 
			\end{figure} 
		
		
		\subsubsection{Programmes COBOL} 
		
			\begin{figure}[h!]
			\caption{Architecture des programmes COBOL (la légende indique quelles sont les parties existantes au début du développement COBOL)}
			\centering
%				\includegraphics[width=12cm]{Analyse_cobol_light.png} 
			\end{figure}
			
			\begin{description}
				 \item[Application web] \hfill \\
				 L'interface web écrite en C\# et Xaml. En réponse aux action des utilisateurs, elle déclenche des requêtes vers le host où résident les programmes COBOL.
				 
				 \item[PSB] \hfill \\
				 Un programme C auto-généré qui fait tous les appels aux sous-programmes COBOL.
				 
				 \item[TWEB] \hfill \\
				 Un programme COBOL qui vérifie que la requête provenant du web est bien formée.
				 
				 \item[Programme métier] \hfill \\
				 Il effectue le routage (appel de sous-programmes logiques selon la requête) et héberge les règles métier (par exemple détermination du bimestre de facturation).
				 
				 \item[Programmes logiques] \hfill \\
				 Une couche de vérification supplémentaire. Les programmes logiques vérifient que le requête vers la base de données est bien formée.
				 
				 \item[Programmes physiques] \hfill \\
				 Exécution de la requête SQL et renvoie des informations.			 			 
			\end{description}
			
			Il est à noter que l'écosystême COBOL ne suit pas le modèle middleware mais le modèle programme appellant- programme apppelé avec stockage des informations de retour dans des variables globales situées dans
			des programmes COBOL dédiés (appelés Clause Copy).

			
\section{Remerciements}

	Plusieurs personnes m’ont apporté une aide significative sur ce projet et je tiens à les remercier chaleureusement ici. 

	\begin{itemize}
		\item Caroline \textsc{SCHIPPERS}, mon maître de stage, qui a réussi malgré son emploi du temps chargé à m’accompagner et me guider tout au long de mon projet
		\item Fabrice \textsc{WENDLING}, analyste développeur, qui m’a beaucoup appris sur le développement web au sein du Framework Devbooster
		\item Orian \textsc{LE SOBRE}, analyste développeur, qui m’a épaulé pour l’analyse et le développement en COBOL, et dans la découverte du système d’information du CM-CIC
		\item Toute l’équipe Gestion de la Facturation, qui m’a cordialement accueilli, aidé et a répondu à toutes mes questions
	\end{itemize}

  
  
\section{Conclusion}


		Ce stage a été, malgré un rythme irrégulier, extrêmement formateur. \\
		
		
		Du point de vue technique, j'ai aperçu le fonctionnement d'un grand systême fonctionnant sur des mainframes avec IBM System z. 
		En outre j'ai discerné les différences d'un site intranet avec un site internet, en particulier au niveau des utilisateurs et de leurs habitudes. Mon application web intranet a donc dû afficher 
		le même design, et les éléments graphiques ont du être agencé exactement comme sur les autres pages de l'intranet, avec la même signification, pour ne pas perturber les utilisateurs, et faciliter leur navigation.
		J'ai aussi beaucoup appris au niveau de l'ergonomie d'un site web, qui est cruciale pour l'expérience utilisateur, surtout dans le contexte d'une application métier qui a vocation a être utilisée
		le plus simplement et rapidement possible. Le but est de gagner en productivité et cela repose en grande partie sur l'ergonomie .\\
		De même des choses nouvelles pour moi ont été les conventions de codage et la structure type d'un programme instaurées par l'entreprise.
		
		
		Au niveau des langages utilisées, le COBOL et ses écueils ont été une totale découverte pour moi, tandis que le C\# a été une semi-découverte puisque je n'en avais jamais fait directement mais 
		j'avais déjà programmé avec des langages ressemblants (C, C++, Java).
		Du point de vue des frameworks, cela a aussi été un apprentissage puisque j'ai utilisé un framework web interne (Devbooster) qui différait du framework web que je connaissais jusque-là: Symfony 2.
		Un autre framework utilisé a été .Net, mais celui-ci était très bien documenté et facile à prendre en main donc celà n'a pas posé de problème.
		
		Un point qui m'a frappé et sur lequel l'accent a été porté durant mon stage est l'architecture, l'analyse et la documentation d'un projet. J'ai en effet eu la chance, en étant guidé, d'élaborer
		l'architecture complète de mon application. L'analyse a aussi été une expérience nouvelle pour moi : lister les contraintes, les fonctionnalités de l'application, les algorithmes utilisés, apporter les solutions adéquates,
		s'adapter aux changements indiqués par les utilisateurs; préparer ainsi en amont le développement a été des plus instructifs.\\
		
		
		Mais la pierre angulaire de ce stage, qui m'a apporté le plus, a été l'enjeu de migration d'une application existante, ce qui ne s'apprend pas à l'école. En effet, j'ai dû trouver le bon équilibre
		entre le contentement des utilisateurs, la conservation des fonctionnalités existantes sans dégradation et l'apport de nouvelles fonctionnalités, sans avoir de documentation tout à fait explicite.
		Certains points ont donc nécessité une recherche dans le code COBOL et SQL existant, notamment l'algorithme de détermination du bimestre de facturation,et la structure de la base de données.
		
		
		D'un point de vue moins technique, c'était mon souhait de faire l'expérience d'une très grande entreprise et découvrir le monde de la banquassurance et ce stage a été parfait sur ces deux points.
		En effet, j'ai expérimenté les contraintes spécifiques à ce domaine (sécurité, fiabilité, robustesse, normes de qualité, délais etc), le fonctionnement d'une grosse entreprise (outils de reporting et monitoring,
		habilitations, etc) et l'informatique de gestion.
		
		En conclusion, l'apport de ce stage est inestimable tant pour mes compétences informatiques que pour mon orientation et mon parcours dans le monde du travail.
	
	

% \appendix
\section{Annexes}
	\subsection{Captures d'écran de l'application existante}
	\newpage
		\begin{figure}[h!]
			\caption{Le menu principal}
			\centering
%				\includegraphics[width=12cm]{ftl_menu.png} 
		\end{figure} 
		
		\begin{figure}[h!]
			\caption{La liste des liaisons louées}
			\centering
%				\includegraphics[width=12cm]{ftl_liste_liaisons.png} 
		\end{figure}
		
		\begin{figure}[h!]
			\caption{La modification d'une liaison}
			\centering
%				\includegraphics[width=12cm]{ftl_modif_liaison.png} 
		\end{figure} 
		
		\begin{figure}[h!]
			\caption{La liste des comptes}
			\centering
%				\includegraphics[width=12cm]{ftl_liste_comptes.png} 
		\end{figure} 

		\begin{figure}[h!]
			\caption{La génération de la facturation}
			\centering
%				\includegraphics[width=12cm]{ftl_facturation.png} 
		\end{figure} 
			
		
		\newpage
		
		
		\subsection{Captures d'écran de l'application web}
		\newpage
		\begin{figure}[h!]
			\caption{La liste des liaisons (écran d'accueil)}
			\centering
%				\includegraphics[width=12cm]{liste_liaisons.png} 
		\end{figure} 
		
		\begin{figure}[h!]
			\caption{La liste des liaisons après une génération de facturation réussie}
			\centering
%				\includegraphics[width=12cm]{liste_liaisons_factu_ok.png} 
		\end{figure} 
		
		\begin{figure}[h!]
			\caption{La liste des liaisons après une modification de liaison réussie}
			\centering
%				\includegraphics[width=12cm]{liste_liaisons_modif_ok.png} 
		\end{figure} 
		
		\begin{figure}[h!]
			\caption{La liste des liaisons après une recherche de liaisons réussie}
			\centering
%				\includegraphics[width=12cm]{liste_liaisons_recherche_ok.png} 
		\end{figure}
		
		\begin{figure}[h!]
			\caption{La liste des liaisons après une suppression de liaison réussie}
			\centering
%				\includegraphics[width=12cm]{liste_liaisons_suppr_ok.png} 
		\end{figure} 

		\begin{figure}[h!]
			\caption{Le détail d'une liaison}
			\centering
%				\includegraphics[width=12cm]{detail_liaison.png} 
		\end{figure}
			

		
		\begin{figure}[h!]
			\caption{La génération de la facturation}
			\centering
%				\includegraphics[width=12cm]{facturation.png} 
		\end{figure} 
		
		\begin{figure}[h!]
			\caption{La modification d'une liaison}
			\centering
%				\includegraphics[width=12cm]{formulaire_liaison.png} 
		\end{figure} 
		
		
		
		\begin{figure}[h!]
			\caption{La suppression d'une liaison}
			\centering
%				\includegraphics[width=12cm]{supprimer_liaison.png} 
		\end{figure}
		
			
		\begin{figure}[h!]
			\caption{La liste des comptes}
			\centering
%				\includegraphics[width=12cm]{liste_comptes.png} 
		\end{figure} 
		
		\begin{figure}[h!]
			\caption{La liste des comptes après une création de compte réussie}
			\centering
%				\includegraphics[width=12cm]{liste_comptes_creat_ok.png} 
		\end{figure}
		
		\begin{figure}[h!]
			\caption{La liste des comptes après une modification de compte réussie}
			\centering
%				\includegraphics[width=12cm]{liste_comptes_modif_ok.png} 
		\end{figure} 
		\clearpage
		\begin{figure}[h!]
			\caption{La liste des comptes après une suppression de compte réussie}
			\centering
%				\includegraphics[width=12cm]{liste_comptes_suppr_ok.png} 
		\end{figure} 
		
		\begin{figure}[h!]
			\caption{La modification d'un compte}
			\centering
%				\includegraphics[width=12cm]{modif_compte.png} 
		\end{figure} 
		
		\begin{figure}[h!]
			\caption{La suppression d'un compte}
			\centering
%				\includegraphics[width=12cm]{supprimer_compte.png} 
		\end{figure} 
	
	
	
	\subsection{Bibliographie}
	
		\begin{description}		
				\item[Cobol] \hfill \\
					\begin{itemize}
						\item IBM Formation technique, Cours DA510F Support élève, Sixième édition Octobre 2002, Référence IBM-ZRF2-7595-P6
						\item IBM Formation technique, Cours DA510F Guide d'exercices, Cinquième édition Octobre 2002, Référence IBM-ZRF2-7596-P6
						\item \url{ibm.com/fr/formation}
						\item \url{ibm.com/training/global}
						\item documentation interne(portail pixis)
					\end{itemize}
					
				\item[C\#] \hfill \\
					\begin{itemize}
						\item \url{www.dotnetperls.com}
						\item \url{msdn.microsoft.com/fr-fr/library}
					\end{itemize}

				\item[Devbooster] \hfill \\
					Documentation interne (portail Pixis).
		\end{description}

		
	\subsection{Glossaire}
		\begin{description} 
			\item[.Net] \hfill \\
				Framework en langage C\# développé par Microsoft, qui fournit notamment des structures de données,
				des interfaces utilisateurs, des accès aux bases de données et des composants pour le développement web.
			\item[Ajax] \hfill \\
				  Asynchronous JavaScript and Xml. Technique de développement web utilisée pour rafraîchir partiellement
				  une page web suite à une action de l'utilisateur.
			\item[C] \hfill \\
				Langage de programmation impératif réputé pour ses performances.
			\item[C\#] \hfill \\
				Langage de programmation orienté objet développé par Microsoft.
			\item[C++] \hfill \\
				Langage de programmation orienté objet et impératif réputé pour ses performances.
			\item[Clause Copy] \hfill \\
				Programme COBOL où sont définies des variables globales utilisées par d'autres programmes COBOL.
				L'échange de données entre ces programmes se fait par la lecture/écriture de ces variables dans la Clause Copy.		
			\item[Cobol] \hfill \\
				 COmmon Business-Oriented Language. Un des premiers langages de programmation, utilisé dans les domaines de la finance, gestion, et de l'administration.
				 Il est utilisé pour édicter les règles métier et effectuer les traitements.
			\item[DB2] \hfill \\
				Base de données relationnelle développée par IBM, accessible au moyen de requêtes SQL.
			\item[Devbooster] \hfill \\
				Framework interne à EI pour développer des applications Web, basé sur le modèle MVC. La Vue est développée en Xaml, le Modèle en Xml et les Contrôleurs en C\#.
			\item[Framework] \hfill \\
				Ensemblede composants (classes, librairies) réutilisables dans d'autres programmes.
			\item[Host]  \hfill \\
				Désigne les mainframes.
			\item[HTML] \hfill \\
				HyperText Markup Language. Langage utilisé dans la création de pages web.
			\item[IHM] \hfill \\
				Interface Homme-Machine. Terme employé pour décrire la partie d'une application qui est visible par l'utilisateur et par laquelle il intéragit avec ladite application.
			\item[Java] \hfill \\
				Langage orienté objet développé par Sun Microsystems.
			\item[Javascript] \hfill \\
				Langage de script notamment utilisé en développement web, côté client mais aussi côté serveur.
			\item[LTNV] \hfill \\
				Format d'échange de données interne à EI.
			\item[Mainframe] \hfill \\
				Abréviation de "Mainframe Computer". Ordinateurs produits par IBM, surtout utilisés par des grosses entreprises et gouvernements,
				ayant des besoins critiques en termes de resources et de robustesse. Ils disposent notamment de composants redondants pour éviter la perte de données et assurer la continuité de l'activité.
			\item[Pixis] \hfill \\
				Portail de documentation situé sur l'intranet du Crédit Mutuel.
			\item[PSB] \hfill \\
				 Un programme C auto-généré qui fait tous les appels aux sous-programmes COBOL. Propre à EI.
			\item[SQL] \hfill \\
				Structured Query Language. Langage utilisé pour la manipulation de données stockées en base de données.
			\item[System z] \hfill \\
				Famille de mainframes proposés par IBM. "z" signifie "zero downtime".
			\item[TWEB] \hfill \\
				Un programme COBOL qui vérifie que la requête provenant du web est bien formée.
			\item[Webfarm] \hfill \\
				Cluster de serveurs destinés à héberger des sites Web et à répondre à un nombre élevé de requêtes.
			\item[Xaml] \hfill \\
				Extensible Application Markup Language. Dialecte Xml développé par Microsoft utilisé pour décrire l'Interface Homme-Machine.
			\item[Xml] \hfill \\
				Extensible Markup Language. Langage descriptif définissant un format d'échange de données compréhensible par l'humain et par la machine.
			\item[z/Os] \hfill \\
				Systême d'exploitation 64 bits produit par IBM, utilisé sur les mainframes.
			\end{description}
		
\end{document}

